\section{Context Free Languages}

\subsection{Context Free Grammars}

  When a person designs a programming language, they need to determine its \textit{syntax}. That is, the designer decides which strings correspond to valid programs, and which ones do not (i.e. which strings contain a syntax error). To ensure that a compiler or interpreter always halts when checking for syntax errors, language designers typically \textit{do not} use a general Turing-complete mechanism to express their syntax. Rather, they use a \textit{restricted} computational model, most often being \textit{context free grammars}. 
	
	\begin{example}[Arithmetic Function]
		Consider the function $ARITH: \Sigma^* \longrightarrow \{0,1\}$ that takes as input a string $x$ over alphabet 
		\[\Sigma = \{(, ), +, -, \times, \div, 0, 1, 2, 3, 4, 5, 6, 7, 8, 9\}\]
		and returns $1$ if and only if the string $x$ represents a valid arithmetic expression. Intuitively, we build expressions by applying an operation such as $+, -, \times, \div$ to smaller expressions or enclosing them in parentheses. More precisely, we can make the following definitions: 
		\begin{enumerate}
			\item A \textit{digit} is one of the symbols $0, 1, 2, 3, 4, 5, 6, 7, 8, 9$. 
			\item A \textit{number} is a sequence of digits (we will drop the condition that the sequence does not have a leading zero)
			\item An \textit{operation} is one of $+, -, \times, \div$. 
			\item An \textit{expression} has either the form 
			\begin{enumerate}
				\item \textit{"number"}
				\item \textit{"sub-expression1 operation sub-expression2}
				\item \textit{"(sub-expression1)"}
			\end{enumerate}
			where "sub-expression1" and "sub-expression2" are themselves expressions. Note that this is a recursive function. 
		\end{enumerate}
		A context free grammar (CFG) is a formal way of specifying such conditions, consisting of a set of rules that tell us how to generate strings from smaller components. 
	\end{example}

  \begin{definition}[Context Free]
		A grammar is \textbf{context free}\footnote{For theorists, context sensitive grammars are important, but for engineers, context free grammars matter much more.} if 
    \begin{equation}
      R \subset V \times (V \cup \Sigma)^\ast
    \end{equation}
    That is, it asserts that the left-hand side can only contain a single non-terminal. A language $L$ is \textbf{context-free} if $L = L(G)$ for some context-free grammar $G$. 
  \end{definition}

  \begin{example}
		The example of well-formed arithmetic expressions can be captured formally by the following context free grammar. 
		\begin{enumerate}
			\item The alphabet $\Sigma$ is $\{(, ), +, -, \times, \div, 0, 1, 2, 3, 4, 5, 6, 7, 8, 9\}$. 
			\item The variables are $V = \{expression, number, digit, operation\}$
			\item The rules are the set $R$ containing the following 19 rules: 
			\begin{enumerate}
				\item 4 Rules: $operation \implies +$, $operation \implies -$, $operation \implies \times$, $operation \implies \div$
				\item 10 Rules: $digit \implies 0$, $digit \implies 1$, ..., $digit \implies 9$
				\item Rule: $number \implies digit$
				\item Rule: $number \implies digit number$
				\item Rule: $expression \implies number$
				\item Rule: $expression \implies expression\; operation\; expression$
				\item Rule: $expression \implies (expression)$
			\end{enumerate}
			\item The starting variable is $expression$.
		\end{enumerate}
  \end{example}

  This allows for a parse-tree structure because each variable ``branches'' into a new string independently of its neighbors. 

  \begin{definition}[Ambiguous]
    A context-free grammar is \textbf{ambiguous} if there is a string that can have more than one valid derivation tree. 
  \end{definition}

  \begin{example}[Ambiguous Context Free Grammar]
    Let $V = \{E\}$ be the non-terminals, and $\{+, -, \texttt{num}\}$ be our terminals, which come from our lexer. Then, our production rules can look something like. 
    \begin{enumerate}
      \item $E \to E + E$. 
      \item $E \to E - E$. 
      \item $E \to \texttt{num}$. 
    \end{enumerate}
    So basically, if we have some word, then we can replace any $E$ in the word by any of the three choices on the right hand side. For example, we can do the following sequence of derivations. 
    \begin{align}
      E & \derives E + E \\ 
        & \derives E - E + E \\
        & \derives E - \texttt{num} + E\\
        & \derives E + E - \texttt{num} + E 
    \end{align}
    You can keep doing this until there is no more production rules you can apply. For context free grammars, you basically do this until there are only terminals left, e.g. $\texttt{num} + \texttt{num} - \texttt{num} + \texttt{num}$. But note that this is ambiguous since there are multiple ways to derive. For example, if we wanted to get something like \texttt{num} - \texttt{num} - \texttt{num}, we can do it in either of the following ways. 
    \begin{align}
      E & \derives E - E \\ 
        & \derives E - (E - E) \\ 
        & \derives \texttt{num} - (\texttt{num} - \texttt{num}) \\
      E & \derives E - E \\ 
        & \derives (E - E) - E \\ 
        & \derives (\texttt{num} - \texttt{num}) - \texttt{num} 
    \end{align}
  \end{example}

  This is not good, so we want to modify our grammar so that it is safe from these ambiguities. 

  \begin{example}[Non-Ambiguous Context-Free Grammar]
    If we have a grammar with the following production rules. 
    \begin{enumerate}
      \item $E \to E + \texttt{num}$
      \item $E \to E - \texttt{num}$ 
      \item $E \to \texttt{num}$
    \end{enumerate}
    Then, we have no ambiguities since there is only one possible way 
    \begin{align}
      E & \derives E - \texttt{num} \\ 
        & \derives (E - \texttt{num}) - \texttt{num} \\
        & \derives (\texttt{num} - \texttt{num}) - \texttt{num} 
    \end{align}
  \end{example}

  Note that figuring out whether a grammar is ambiguous or unambiguous is a bit tricky. It is in fact undecidable. 

